\documentclass[10pt,a4paper]{amsart}
\usepackage[margin=19mm]{geometry}
\usepackage{amsmath,amssymb,mathtools}
\usepackage[scr=rsfs,cal=euler]{mathalfa}
\usepackage{xfrac}
\usepackage[unicode,hidelinks,bookmarks=false]{hyperref}
\newcommand{\ie}{i.e.\ }
\newcommand{\cf}{cf.\ }
\newcommand{\resp}{respectively\ }
\newcommand{\Subsection}{\subsection}
\providecommand{\nobreakdash}{\nobreak\hbox{-}}
\setlength{\emergencystretch}{2em}
\multlinegap0pt
\usepackage{bm}
\usepackage{bbm}
\usepackage{dsfont}
\usepackage{xspace}
\usepackage{cancel}
\usepackage{mathtools}
\usepackage{amsthm}
\makeatletter
\@ifundefined{thm}{\newtheorem{thm}{Thm}}{}
\@ifundefined{prop}{\newtheorem{prop}[thm]{Proposition}}{}
\makeatother
\newcommand{\Bea}{\begin{eqnarray*}}
\newcommand{\Eea}{\end{eqnarray*}}
\title[Metrics on $C^{\ast}$-algebras of \'Etale groupoids from len]{Metrics on $C^{\ast}$-algebras of \'Etale groupoids from length functions}
\author{Arnab Chattopadhyay, Md Amir Hossain,, Soumalya Joardar}
\date{Source: arXiv:2504.13530v1 (CC BY 4.0)}
\begin{document}
\maketitle

\noindent\emph{Source and licence.} Metrics on $C^{\ast}$-algebras of \'Etale groupoids from length functions. Version arXiv:2504.13530v1 (CC BY 4.0), distributed under the Creative Commons Attribution 4.0 International licence, as recorded in the arXiv licence metadata for this exact version and shown on the abstract page https://arxiv.org/abs/2504.13530v1. Full rights notice: licenses/arxiv-2504.13530-CC-BY-4.0.txt.\par\smallskip

\noindent\textit{Editorial scope.} Counted unit: the Prop whose source label is \texttt{prop21} in the article (source0.tex lines 404--408, source label \texttt{prop21}), delivered together with the article's own complete original proof of it (source0.tex lines 410--472, one \texttt{proof} environment of 63 lines). No mathematics is written, repaired or re-derived: statement and proof are the article's own text line for line. Required context retained uncounted: none. Every definition, display and label of those lines is retained unchanged. Omitted from this package and not redistributed: 1--403, 409--409, 473--824. Nothing inside a retained excerpt is omitted or altered: every retained body is delivered exactly as the article's lines are written, so no omission receipt of any kind accompanies this package. Citations: the retained excerpts carry no citation command at all, so no bibliography entry of the article is transcribed and no reference is redistributed. Reference adaptation: none was needed and none was made; every reference inside the delivered unit resolves inside it, so no unresolved reference or citation is produced. Printed slips kept literal: no printed word, symbol, hypothesis or number of the counted statement or of its proof was altered. Layout and class: the article's own class is \texttt{amsart} with packages including amssymb, epsfig, amsfonts, amsmath, euscript, amscd, amsthm, enumitem, natbib, color, mathtools, hyperref, marginnote, ulem; the wrapper is the lane's pinned \texttt{amsart} editorial wrapper at 10pt on A4, which restates the theorem-like environments the retained text needs and adds the article's own macro definitions that the retained text needs (\texttt{\textbackslash Bea}, \texttt{\textbackslash Eea}), with extra wrapper packages (bm, bbm, dsfont, xspace, cancel, mathtools). The delivered numbering is the unit's own. Assets: no retained span contains a figure environment, a graphics-inclusion command, a PDF-page inclusion, a table or a TikZ picture; no image file is copied into this package, the package declares no asset dependency and no image file is redistributed with it. Licence: version arXiv:2504.13530v1 of this article is distributed under the Creative Commons Attribution 4.0 International licence, as recorded in the arXiv licence metadata for this exact version and shown on the abstract page https://arxiv.org/abs/2504.13530v1. A full rights notice is delivered at licenses/arxiv-2504.13530-CC-BY-4.0.txt. Characters: every retained excerpt is pure ASCII, so the delivered engine, XeLaTeX with its default Latin Modern text font, reports no missing character.
			\begin{prop}\label{prop21}
				For $f \in C_c (\mathcal G),$ the operator $\Delta^k (f)$ is adjointable and 
				$$\Delta^k (f)^{\ast} = (-1)^k \Delta^k (f^{\ast}).$$ 
				%where $f^{\ast} (\gamma) = \overline {f \left (\gamma^{-1} \right )},$ for all $\gamma \in \mathcal G.$
			\end{prop}

			\begin{proof}
				Since $C_c (\mathcal G)$ is dense in $L^2 (\mathcal G)$ and $\Delta^k (f)$ is bounded, in order to show the desired equality it is enough to show that for any $\xi, \eta \in C_c (\mathcal G),$ we have 
				
				$$\bigl \langle \bigl \langle \Delta^k (f) (\xi), \eta \bigr \rangle \bigr \rangle = (-1)^k \bigl \langle \bigl \langle \xi, \Delta^k (f^{\ast}) (\eta) \bigr \rangle \bigr \rangle.$$
				
				We prove the above equality for the case $k=1$. For the general case a regular induction argument will suffice and we leave the induction argument to the reader. To that end, for $f, \xi, \eta \in C_c (\mathcal G)$ and $x \in \mathcal G^{(0)}$ we have
				\Bea
				\left \langle \left \langle M_{\ell} (\lambda (f) (\xi)), \eta \right \rangle \right \rangle (x)
				= \sum\limits_{\gamma \in \mathcal G_x} \ell (\gamma) \overline {\lambda (f) (\xi) (\gamma)}\ \eta (\gamma)  & = &  \sum\limits_{\gamma \in \mathcal G_x} \ell (\gamma) \sum\limits_{\beta \in \mathcal G_x} \overline {f \left (\gamma \beta^{-1} \right )}\ \overline {\xi (\beta)}\ \eta (\gamma) \\  & = &  \sum\limits_{\gamma \in \mathcal G_x} \sum\limits_{\beta \in G_x} \ell (\gamma) \overline {f \left (\gamma \beta^{-1} \right )}\ \overline {\xi (\beta)}\ \eta (\gamma).
				\Eea
				Using Fubini's theorem the last quantity can be written as
				\Bea
				\sum\limits_{\beta \in \mathcal G_x} \sum\limits_{\gamma \in \mathcal G_x} \ell (\gamma) \overline {f \left (\gamma \beta^{-1} \right )}\ \overline {\xi (\beta)}\ \eta (\gamma)  & = &  \sum\limits_{\beta \in \mathcal G_x} \overline {\xi (\beta)}\ \sum\limits_{\gamma \in \mathcal G_x} \ell (\gamma) f^{\ast} \left (\beta \gamma^{-1} \right )\ \eta (\gamma) \\  & = &  \sum\limits_{\beta \in \mathcal G_x} \overline {\xi (\beta)}\ \sum\limits_{\gamma \in \mathcal G_x} f^{\ast} \left (\beta \gamma^{-1} \right )\ M_{\ell} (\eta) (\gamma) \\ & = &  \sum\limits_{\beta \in \mathcal G_x} \overline {\xi (\beta)} \left (f^{\ast} \ast M_{\ell} (\eta) \right ) (\beta) \\ & = & \sum\limits_{\beta \in \mathcal G_x} \overline {\xi (\beta)} \lambda (f^{\ast}) \left (M_{\ell} (\eta) \right ) (\beta) \\ & = & \left \langle \left \langle \xi, \lambda (f^{\ast}) \left (M_{\ell} (\eta) \right ) \right \rangle \right \rangle (x).  
				\Eea
				Again for \(f, \xi, \eta \in C_c (\mathcal G)\) and \(x\in \mathcal{G}^{(0)}\), we have
				\Bea
					\left \langle \left \langle \lambda (f) \left (M_{\ell} (\xi) \right ), \eta \right \rangle \right \rangle (x) =\sum\limits_{\gamma \in \mathcal G_x} \overline {\lambda (f) \left (M_{\ell} (\xi) \right ) (\gamma)}\ \eta (\gamma) & = & \sum\limits_{\gamma \in \mathcal G_x} \overline {\left (f \ast M_{\ell} (\xi) \right ) (\gamma)}\ \eta (\gamma) \\ & = &  \sum\limits_{\gamma \in \mathcal G_x} \sum\limits_{\beta \in \mathcal G_x} \overline {f \left (\gamma \beta^{-1} \right )}\ \overline {M_{\ell} (\xi) (\beta)}\ \eta (\gamma).
				\Eea
				Again, using Fubini's theorem the last term becomes
				\Bea
				\sum\limits_{\beta \in \mathcal G_x} \sum\limits_{\gamma \in \mathcal G_x} \overline {f \left (\gamma \beta^{-1} \right )}\ \ell (\beta)\ \overline {\xi (\beta)}\ \eta (\gamma)  & = &  \sum\limits_{\beta \in \mathcal G_x} \ell (\beta) \overline {\xi (\beta)} \sum\limits_{\gamma \in \mathcal G_x} \overline {f \left (\gamma \beta^{-1} \right )}\ \eta (\gamma) \\ & = &  \sum\limits_{\beta \in \mathcal G_x} \ell (\beta) \overline {\xi (\beta)} \sum\limits_{\gamma \in \mathcal G_x} f^{\ast} \left (\beta \gamma^{-1} \right )\ \eta (\gamma) \\ & = &  \sum\limits_{\beta \in \mathcal G_x} \ell (\beta) \overline {\xi (\beta)} \left (f^{\ast} \ast \eta \right ) (\beta) \\ & = &  \sum\limits_{\beta \in \mathcal G_x} \overline {\xi (\beta)} M_{\ell} \left (f^{\ast} \ast \eta \right ) (\beta) \\ & = & \left \langle \left \langle \xi, M_{\ell} \left (\lambda (f^{\ast}) (\eta) \right ) \right \rangle \right \rangle (x).
				\Eea
				
%				\Bea 
%				&&\left \langle \left \langle M_{\ell} (\lambda (f) (\xi)), \eta \right \rangle \right \rangle (x)\\
%				& = & \sum\limits_{\gamma \in \mathcal G_x} \ell (\gamma) \overline {\lambda (f) (\xi) (\gamma)}\ \eta (\gamma) \\ & = & \sum\limits_{\gamma \in \mathcal G_x} \ell (\gamma) \sum\limits_{\beta \in \mathcal G_x} \overline {f \left (\gamma \beta^{-1} \right )}\ \overline {\xi (\beta)}\ \eta (\gamma) \\ & = & \sum\limits_{\gamma \in \mathcal G_x} \sum\limits_{\beta \in G_x} \ell (\gamma) \overline {f \left (\gamma \beta^{-1} \right )}\ \overline {\xi (\beta)}\ \eta (\gamma) \\ & \stackrel{\text {Fubini}}{=} & \sum\limits_{\beta \in \mathcal G_x} \sum\limits_{\gamma \in \mathcal G_x} \ell (\gamma) \overline {f \left (\gamma \beta^{-1} \right )}\ \overline {\xi (\beta)}\ \eta (\gamma) \\ & = & \sum\limits_{\beta \in \mathcal G_x} \overline {\xi (\beta)}\ \sum\limits_{\gamma \in \mathcal G_x} \ell (\gamma) f^{\ast} \left (\beta \gamma^{-1} \right )\ \eta (\gamma) \\ & = & \sum\limits_{\beta \in \mathcal G_x} \overline {\xi (\beta)}\ \sum\limits_{\gamma \in \mathcal G_x} f^{\ast} \left (\beta \gamma^{-1} \right )\ M_{\ell} (\eta) (\gamma) \\ & = & \sum\limits_{\beta \in \mathcal G_x} \overline {\xi (\beta)} \left (f^{\ast} \ast M_{\ell} (\eta) \right ) (\beta) \\ & = & \sum\limits_{\beta \in \mathcal G_x} \overline {\xi (\beta)} \lambda (f^{\ast}) \left (M_{\ell} (\eta) \right ) (\beta) \\ & = & \left \langle \left \langle \xi, \lambda (f^{\ast}) \left (M_{\ell} (\eta) \right ) \right \rangle \right \rangle (x) \  
%				\Eea
				
%				and, 
				
%				\Bea
%				&&\left \langle \left \langle \lambda (f) \left (M_{\ell} (\xi) \right ), \eta \right \rangle \right \rangle (x)\\
%				& = & \sum\limits_{\gamma \in \mathcal G_x} \overline {\lambda (f) \left (M_{\ell} (\xi) \right ) (\gamma)}\ \eta (\gamma) \\ & = & \sum\limits_{\gamma \in \mathcal G_x} \overline {\left (f \ast M_{\ell} (\xi) \right ) (\gamma)}\ \eta (\gamma) \\ & = & \sum\limits_{\gamma \in \mathcal G_x} \sum\limits_{\beta \in \mathcal G_x} \overline {f \left (\gamma \beta^{-1} \right )}\ \overline {M_{\ell} (\xi) (\beta)}\ \eta (\gamma) \\ & \stackrel{\text {Fubini}} {=} & \sum\limits_{\beta \in \mathcal G_x} \sum\limits_{\gamma \in \mathcal G_x} \overline {f \left (\gamma \beta^{-1} \right )}\ \ell (\beta)\ \overline {\xi (\beta)}\ \eta (\gamma) \\ & = & \sum\limits_{\beta \in \mathcal G_x} \ell (\beta) \overline {\xi (\beta)} \sum\limits_{\gamma \in \mathcal G_x} \overline {f \left (\gamma \beta^{-1} \right )}\ \eta (\gamma) \\ & = & \sum\limits_{\beta \in \mathcal G_x} \ell (\beta) \overline {\xi (\beta)} \sum\limits_{\gamma \in \mathcal G_x} f^{\ast} \left (\beta \gamma^{-1} \right )\ \eta (\gamma) \\ & = & \sum\limits_{\beta \in \mathcal G_x} \ell (\beta) \overline {\xi (\beta)} \left (f^{\ast} \ast \eta \right ) (\beta) \\ & = & \sum\limits_{\beta \in \mathcal G_x} \overline {\xi (\beta)} M_{\ell} \left (f^{\ast} \ast \eta \right ) (\beta) \\ & = & \left \langle \left \langle \xi, M_{\ell} \left (\lambda (f^{\ast}) (\eta) \right ) \right \rangle \right \rangle (x)
%				\Eea
				
				
Therefore, we have 
\Bea
	\left \langle \left \langle \Delta (f) (\xi), \eta \right \rangle \right \rangle & = &  \left \langle \left \langle M_{\ell} (\lambda (f) (\xi)), \eta \right \rangle \right \rangle - \left \langle \left \langle \lambda (f) (M_{\ell} (\xi)), \eta \right \rangle \right \rangle \\ & = &  \left \langle \left \langle \xi, \lambda (f^{\ast}) \left (M_{\ell} (\eta) \right ) \right \rangle \right \rangle - \left \langle \left \langle \xi, M_{\ell} \left (\lambda (f^{\ast}) (\eta) \right ) \right \rangle \right \rangle \\ & = &  - \left \langle \left \langle \xi, \Delta (f^{\ast}) (\eta) \right \rangle \right \rangle.  
\Eea
				
%				\Bea 
%				\left \langle \left \langle \Delta (f) (\xi), \eta \right \rangle \right \rangle & = & \left \langle \left \langle M_{\ell} (\lambda (f) (\xi)), \eta \right \rangle \right \rangle - \left \langle \left \langle \lambda (f) (M_{\ell} (\xi)), \eta \right \rangle \right \rangle \\ & = & \left \langle \left \langle \xi, \lambda (f^{\ast}) \left (M_{\ell} (\eta) \right ) \right \rangle \right \rangle - \left \langle \left \langle \xi, M_{\ell} \left (\lambda (f^{\ast}) (\eta) \right ) \right \rangle \right \rangle \\ & = & - \left \langle \left \langle \xi, \Delta (f^{\ast}) (\eta) \right \rangle \right \rangle.  
%				\Eea
		This shows that $\Delta (f)^{\ast} = - \Delta (f^{\ast})$ for any $f \in C_c (\mathcal G),$  proving the equality for $k = 1.$ 
				%Let us assume that the equality holds for $k -1$ i.e. for all $f \in C_c (\mathcal G)$ 
				%$$\Delta^{k - 1} (f)^{\ast} = (-1)^{k - 1} \Delta^{k - 1} (f^{\ast}).$$
				%Now using induction hypothesis we have
				%\Bea
				%\left \langle \left \langle M_{\ell} \left (\Delta^{k - 1} (f) (\xi) \right ), \eta \right \rangle \right \rangle (x) & = & \sum\limits_{\gamma \in \mathcal G_x} \overline {M_{\ell} \left (\Delta^{k - 1} (f) (\xi) \right ) (\gamma)}\ \eta (\gamma) \\ & = & \sum\limits_{\gamma \in \mathcal G_x} \ell (\gamma)\ \overline {\Delta^{k - 1} (f) (\xi) (\gamma)}\ \eta (\gamma) \\ & = & \sum\limits_{\gamma \in \mathcal G_x} \overline {\Delta^{k - 1} (f) (\xi) (\gamma)}\ M_{\ell} (\eta) (\gamma) \\ & = & \left \langle \left \langle \Delta^{k - 1} (f) (\xi), M_{\ell} (\eta) \right \rangle \right \rangle (x) \\ & = & \left \langle \left \langle \xi, \Delta^{k-1}(f)^{\ast} \left (M_{\ell} (\eta) \right ) \right \rangle \right \rangle (x) \\ & = & (-1)^{k - 1} \left \langle \left \langle \xi, \Delta^{k - 1} (f^{\ast}) \left (M_{\ell} (\eta) \right ) \right \rangle \right \rangle (x)
				%\Eea
				%and,
				%\Bea
				%\left \langle \left \langle \Delta^{k - 1} (f) \left (M_{\ell} (\xi) \right ), \eta \right \rangle \right \rangle (x) & = & \left \langle \left \langle M_{\ell} (\xi), \Delta^{k - 1}(f)^{\ast} (\eta) \right \rangle \right \rangle (x) \\ & = & (-1)^{k - 1} \left \langle \left \langle M_{\ell} (\xi), \Delta^{k - 1}(f^{\ast}) (\eta) \right \rangle \right \rangle (x) \\ & = & (-1)^{k - 1} \sum\limits_{\gamma \in \mathcal G_x} \overline {M_{\ell} (\xi) (\gamma)}\ \Delta^{k - 1} (f^{\ast}) (\eta) (\gamma) \\ & = & (-1)^{k - 1} \sum\limits_{\gamma \in \mathcal G_x} \ell (\gamma)\ \overline {\xi (\gamma)}\ \Delta^{k - 1} (f^{\ast}) (\eta) (\gamma) \\ & = & (-1)^{k - 1} \sum\limits_{\gamma \in \mathcal G_x} \overline {\xi (\gamma)}\ M_{\ell} \left (\Delta^{k - 1} (f^{\ast}) (\eta) \right ) (\gamma) \\ & = & (-1)^{k - 1} \left \langle \left \langle \xi, M_{\ell} \left (\Delta^{k - 1} (f^{\ast}) (\eta) \right ) \right \rangle \right \rangle (x)    
				%\Eea
				%Therefore
			 %\Bea 
				%&& \left \langle \left \langle \Delta^{k} (f) (\xi), \eta \right \rangle \right \rangle \\ & = & \left \langle \left \langle M_{\ell} \left (\Delta^{k - 1} (f) (\xi) \right ), \eta \right \rangle \right \rangle - \left \langle \left \langle \Delta^{k - 1} (f) (M_{\ell} (\xi)), \eta \right \rangle \right \rangle \\ & = & (-1)^{k -1} \left \langle \left \langle \xi, \Delta^{k - 1} (f^{\ast}) \left (M_{\ell} (\eta) \right ) \right \rangle \right \rangle - (-1)^{k - 1} \left \langle \left \langle \xi, M_{\ell} \left (\Delta^{k - 1} (f^{\ast}) (\eta) \right ) \right \rangle \right \rangle \\ & = & (-1)^k \left \langle \left \langle \xi, \Delta^{k} (f^{\ast}) (\eta) \right \rangle \right \rangle   
				%\Eea
				%With a simple induction argument it can be proved that
				%$$\Delta^k (f)^{\ast} = (-1)^{k} \Delta^k (f^{\ast}),$$
				%for all $f \in C_c (\mathcal G),$ as required. This completes the proof.
			\end{proof}

\end{document}
